\clearpage
\VerdierPhysicalPageAnchor{0134}%
\begingroup
\small
\setlength{\emergencystretch}{1.5em}%
\setlength{\abovedisplayskip}{1pt plus 1pt minus 1pt}%
\setlength{\belowdisplayskip}{1pt plus 1pt minus 1pt}%

\VerdierSemanticAnchor{ch2-2-1-15}
\noindent\textbf{2.1.15.}\quad
\VerdierSemanticAnchor{ch2-2-1-15-associated-system}%
设 $\mathcal D'$ 为三角范畴 $\mathcal D$ 的满三角子范畴。乘法系
$S(\mathcal D')$ (2.1.7) 称为\emph{与子范畴 $\mathcal D'$ 关联的
乘法系}。

\VerdierSemanticAnchor{ch2-2-1-16}
\smallskip
\noindent\textbf{命题 2.1.16.}\quad
\begingroup\itshape
\VerdierSemanticAnchor{ch2-2-1-16-setup}%
设 $F\colon\mathcal D_1\to\mathcal D_2$ 为两个三角范畴之间的正合
函子 (1.1.3)。以 $S(F)$ 表示 $\mathcal D_1$ 中被 $F$ 送到
$\mathcal D_2$ 的同构的态射集合。
\VerdierSemanticAnchor{ch2-2-1-16-kernel-subcategory}%
以 $\mathcal D'(F)$ 表示 $\mathcal D_1$ 的满子范畴，其对象是被
$F$ 送到 $\mathcal D_2$ 的零对象的那些 $\mathcal D_1$ 的对象。

\VerdierSemanticAnchor{ch2-2-1-16-a}
\par\noindent\hspace*{2em}a)\enspace 子范畴 $\mathcal D'(F)$ 是严格满
且饱和的三角子范畴（$\mathcal D_1$ 中每个同构于 $\mathcal D'(F)$
的某个对象的对象，都是 $\mathcal D'(F)$ 的对象）。

\VerdierSemanticAnchor{ch2-2-1-16-b}
\par\noindent\hspace*{2em}b)\enspace 集合 $S(F)$ 是 $\mathcal D_1$
的饱和乘法系。

\VerdierSemanticAnchor{ch2-2-1-16-c}
\par\noindent\hspace*{2em}c)\enspace 集合 $S(F)$ 是与子范畴
$\mathcal D'(F)$ 关联的乘法系 (2.1.15)。
\par\endgroup

\VerdierSemanticAnchor{ch2-2-1-16-proof-omitted}
\noindent 证明很容易，留给读者；读者也可用同样的方法证明下列命题：

\VerdierSemanticAnchor{ch2-2-1-17}
\smallskip
\noindent\textbf{命题 2.1.17.}\quad
\begingroup\itshape
\VerdierSemanticAnchor{ch2-2-1-17-setup}%
设 $H\colon\mathcal D\to\mathcal A$ 为从三角范畴到阿贝尔范畴的
上同调函子。以 $S(H)$ 表示 $\mathcal D$ 中连同其所有平移的像都为
$\mathcal A$ 的同构的态射集合。
\VerdierSemanticAnchor{ch2-2-1-17-kernel-subcategory}%
以 $\mathcal D'(H)$ 表示 $\mathcal D$ 的满子范畴，其对象是连同
其所有平移的像都为 $\mathcal A$ 的零对象的那些
$\mathcal D$ 的对象。

\VerdierSemanticAnchor{ch2-2-1-17-a}
\par\noindent\hspace*{2em}a)\enspace 子范畴 $\mathcal D'(H)$ 是严格满
且饱和的三角子范畴。

\VerdierSemanticAnchor{ch2-2-1-17-b}
\par\noindent\hspace*{2em}b)\enspace 集合 $S(H)$ 是 $\mathcal D$
的饱和乘法系。

\VerdierSemanticAnchor{ch2-2-1-17-c}
\par\noindent\hspace*{2em}c)\enspace 集合 $S(H)$ 是与子范畴
$\mathcal D'(H)$ 关联的乘法系 (2.1.15)。
\par\endgroup

\VerdierSemanticAnchor{ch2-2-2}
\medskip
\noindent\textbf{2.2.\quad 局部化范畴的构造。}

\VerdierSemanticAnchor{ch2-2-2-1}
\smallskip
\noindent\textbf{2.2.1.}\quad
\VerdierSemanticAnchor{ch2-2-2-1-slice-categories}%
设 $\mathcal C$ 为一个范畴（不必加性），$S$ 为 $\mathcal C$ 的乘法系。
对于 $\mathcal C$ 的每个对象 $X$，以 $S/X$（分别地，$X\backslash S$）
表示如下范畴：其对象为 $S$ 中靶为 $X$（分别地，源为 $X$）的态射，
\VerdierSemanticAnchor{ch2-2-2-1-morphism-diagrams}%
其态射为下列交换图：

\VerdierSemanticAnchor{ch2-2-2-1-s-over-x-diagram}
\begin{center}
  \begin{tikzpicture}[>=stealth,line width=.38pt,scale=.90,transform shape]
    \node (Y)  at (-1.55,1.05) {$Y$};
    \node (Yp) at ( 1.55,1.05) {$Y'$};
    \node (X)  at ( 0,-.95) {$X$};
    \draw[->] (Y) -- node[above=1pt] {$u$} (Yp);
    \draw[->] (Y) -- node[left=2pt] {$s$} (X);
    \draw[->] (Yp) -- node[right=2pt] {$t$} (X);
    \node at (3.15,.05) {$s,t\in S$};
  \end{tikzpicture}
  \quad ,
\end{center}

\VerdierSemanticAnchor{ch2-2-2-1-p135-boundary}
\endgroup
